\HLASourcePage{354}{353}
\HLAServiceBookmark{12.1}{1}{Overview}
\HLAServiceBookmark{12.2}{1}{Designators}
\thispagestyle{hla-body}
\begin{hlafixedpage}
  \HLAFixedBlock{90.000}{69.120}{514.000}{\rmfamily\raggedright}{12.00}{13.44}{%
    \textbf{12. Programming language mappings}%
  }
  \HLAFixedBlock{90.000}{101.363}{514.000}{\rmfamily\raggedright}{10.98}{12.30}{%
    \textbf{12.1 Overview}%
  }
  \HLAFixedBlock{90.000}{126.786}{514.000}{\rmfamily\raggedright}{9.96}{11.16}{%
    The services described in this document to this point have been described in a programming-language-%
  }
  \HLAFixedBlock{90.000}{138.785}{514.000}{\rmfamily\raggedright}{9.96}{11.16}{%
    neutral way. This approach was taken because support is provided for multiple programming languages as%
  }
  \HLAFixedBlock{90.000}{150.785}{514.000}{\rmfamily\raggedright}{9.96}{11.16}{%
    described in the annexes to this document. This clause (i.e., Clause 12) provides the reader with a mapping%
  }
  \HLAFixedBlock{90.000}{162.785}{514.000}{\rmfamily\raggedright}{9.96}{11.16}{%
    from the abstract view from the body of this document (Clause 4 through Clause 11) to the programming-%
  }
  \HLAFixedBlock{90.000}{174.785}{514.000}{\rmfamily\raggedright}{9.96}{11.16}{%
    language-specific views in the APIs (Annex B through Annex D). It shall be transparent to the federation%
  }
  \HLAFixedBlock{89.999}{186.785}{514.001}{\rmfamily\raggedright}{9.96}{11.16}{%
    and participating federates which API or implementation language a specific federate uses.%
  }
  \HLAFixedBlock{89.999}{210.784}{514.001}{\rmfamily\raggedright}{9.96}{11.16}{%
    The mappings described in this clause are Ann%
  }
  \HLAFixedBlock{100.019}{234.784}{503.981}{\rmfamily\raggedright}{9.96}{11.16}{%
    ---%
  }
  \HLAFixedBlock{121.979}{234.784}{482.021}{\rmfamily\raggedright}{9.96}{11.16}{%
    Of designators from Clause 4 through Clause 10 to handles and names in the APIs%
  }
  \HLAFixedBlock{100.019}{246.784}{503.981}{\rmfamily\raggedright}{9.96}{11.16}{%
    ---%
  }
  \HLAFixedBlock{121.979}{246.784}{482.021}{\rmfamily\raggedright}{9.96}{11.16}{%
    From services in Clause 4 through Clause 10 to methods in the APIs%
  }
  \HLAFixedBlock{100.019}{258.784}{503.981}{\rmfamily\raggedright}{9.96}{11.16}{%
    ---%
  }
  \HLAFixedBlock{121.979}{258.784}{482.021}{\rmfamily\raggedright}{9.96}{11.16}{%
    From abstract data types in Clause 4 through Clause 10 to classes or data types in the APIs%
  }
  \HLAFixedBlock{90.000}{288.383}{514.000}{\rmfamily\raggedright}{10.98}{12.30}{%
    \textbf{12.2 Designators}%
  }
  \HLAFixedBlock{90.000}{313.806}{514.000}{\rmfamily\raggedright}{9.96}{11.16}{%
    As explained in 1.5, designators are a generalization of both names and handles. In most cases, what is%
  }
  \HLAFixedBlock{90.001}{325.805}{513.999}{\rmfamily\raggedright}{9.96}{11.16}{%
    referred to as a designator in Clause 4 through Clause 10 has both a name and a handle. Which specific form%
  }
  \HLAFixedBlock{90.001}{337.805}{513.999}{\rmfamily\raggedright}{9.96}{11.16}{%
    of a designator is used in the APIs is dependent on the language and service in which it is used. Examine the%
  }
  \HLAFixedBlock{90.001}{349.805}{513.999}{\rmfamily\raggedright}{9.96}{11.16}{%
    APIs to see which form is needed, although in most cases, the handle is used. Clause 10 presents services%
  }
  \HLAFixedBlock{90.001}{361.805}{513.999}{\rmfamily\raggedright}{9.96}{11.16}{%
    that shall be used to map names to handles and vice versa.%
  }
  \HLAFixedBlock{90.001}{385.805}{513.999}{\rmfamily\raggedright}{9.96}{11.16}{%
    Names in most cases (e.g., object class names, attribute names, interaction class names) come from the%
  }
  \HLAFixedBlock{90.001}{397.804}{513.999}{\rmfamily\raggedright}{9.96}{11.16}{%
    FDD. As such, they are known ahead of time to all federates. Object instance names, however, are different.%
  }
  \HLAFixedBlock{90.001}{409.804}{513.999}{\rmfamily\raggedright}{9.96}{11.16}{%
    These names may be chosen by the registering federate, or they may be created dynamically by the RTI. As%
  }
  \HLAFixedBlock{90.001}{421.804}{513.999}{\rmfamily\raggedright}{9.96}{11.16}{%
    such, these names often are not predictable.%
  }
  \HLAFixedBlock{90.001}{445.804}{513.999}{\rmfamily\raggedright}{9.96}{11.16}{%
    Handles in all cases are generated by the RTI. In most cases, where the handles correspond to names from%
  }
  \HLAFixedBlock{90.001}{457.803}{513.999}{\rmfamily\raggedright}{9.96}{11.16}{%
    the FDD, these handles shall be generated in a repeatable manner. In other words, if two federation%
  }
  \HLAFixedBlock{90.001}{469.803}{513.999}{\rmfamily\raggedright}{9.96}{11.16}{%
    executions are created on the same version of the same RTI with the same FDD file, then the handles%
  }
  \HLAFixedBlock{90.001}{481.803}{513.999}{\rmfamily\raggedright}{9.96}{11.16}{%
    assigned to the same name shall be the same in each federation execution. This property facilitates saves and%
  }
  \HLAFixedBlock{90.001}{493.803}{513.999}{\rmfamily\raggedright}{9.96}{11.16}{%
    restores so that handles, rather than just the names, can be saved by federates. This property does not,%
  }
  \HLAFixedBlock{90.001}{505.803}{513.999}{\rmfamily\raggedright}{9.96}{11.16}{%
    however, imply that an RTI assign handles in a manner known a priori to federates.%
  }
  \HLAFixedBlock{90.001}{529.802}{513.999}{\rmfamily\raggedright}{9.96}{11.16}{%
    Handles shall be unique in a federation execution. Each federate shall know the same item by the same%
  }
  \HLAFixedBlock{90.001}{541.802}{513.999}{\rmfamily\raggedright}{9.96}{11.16}{%
    handle. In other words, the federate registering an object instance shall know it by the same handle as%
  }
  \HLAFixedBlock{90.001}{553.802}{513.999}{\rmfamily\raggedright}{9.96}{11.16}{%
    another federate that discovers the object instance.%
  }
  \HLAFixedBlock{90.001}{577.801}{513.999}{\rmfamily\raggedright}{9.96}{11.16}{%
    Often during a federation execution, it is useful for federates to send information identifying an item to%
  }
  \HLAFixedBlock{90.001}{589.801}{513.999}{\rmfamily\raggedright}{9.96}{11.16}{%
    another federate. For example, one federate may wish to send an interaction to another federate informing it%
  }
  \HLAFixedBlock{90.001}{601.801}{513.999}{\rmfamily\raggedright}{9.96}{11.16}{%
    that one particular object instance is of particular importance. To send information identifying an item to%
  }
  \HLAFixedBlock{90.001}{613.801}{513.999}{\rmfamily\raggedright}{9.96}{11.16}{%
    another federate, there are two options: send the name of the item, or send the handle of the item.%
  }
  \HLAFixedBlock{90.001}{637.800}{513.999}{\rmfamily\raggedright}{9.96}{11.16}{%
    To send a handle as an attribute value or as a parameter value, the handle shall be encoded prior to sending,%
  }
  \HLAFixedBlock{90.001}{649.800}{513.999}{\rmfamily\raggedright}{9.96}{11.16}{%
    and it shall be reconstituted by the receiver. The precise means of encoding and reconstituting a handle is%
  }
  \HLAFixedBlock{90.001}{661.800}{513.999}{\rmfamily\raggedright}{9.96}{11.16}{%
    described for each API in 12.11.6.4, 12.12.6.4, and 12.13.3.3. However, a handle encoded using any API%
  }
  \HLAFixedBlock{90.000}{673.800}{514.000}{\rmfamily\raggedright}{9.96}{11.16}{%
    (e.g., C++) shall be reconstitutable in all APIs. Also, a handle of a given type (e.g., an object instance%
  }
  \HLAFixedBlock{90.000}{685.800}{514.000}{\rmfamily\raggedright}{9.96}{11.16}{%
    handle) shall be reconstitutable only as that type.%
  }
\end{hlafixedpage}
\clearpage
